\chapter{CPU and System Bus}

\section{Introduction}

Now that our motherboard is powered and reset and clock signals are available, it is time to connect
the CPU to the rest of the system. 

As introduced in Chapter 2, the CPU is the most important part of the computer, as it is the one
that executes the instructions that make the computer work. The CPU is connected to the rest of the
system through a bus, which is a set of wires that allow the CPU to communicate with the other
components of the system. 

In this chapter, we will see how the CPU is connected to the rest of the system and how it
communicates with the rest of components through the system bus.

\section{Circuit Description}

The CPU used in the MSX System is the Zilog Z80. This CPU is connected to the rest of the system
through the system bus, as shown in the schematic diagram in Figure
\ref{fig:artemisa-schematic-cpu}.

\begin{figure}[htbp]
    \centering
    \includegraphics[width=\linewidth,trim={1cm 10.5cm 2cm 8cm},clip]{figures/artemisa-schematic-cpu}
    \caption{Schematic diagram of CPU and its interface to the system bus}
    \label{fig:artemisa-schematic-cpu}
\end{figure}

In essence, the CPU is connected to the system bus through a set signals that may be classified into
three groups:

\begin{itemize}
    \item {\bf Address bus}: This group of signals is used to send the address of the memory
    location or the port of the I/O device that the CPU wants to read or write. We can see the
    address bus is fully controlled by the CPU, and they are considered output signals from the CPU
    perspective at the top right of Figure \ref{fig:artemisa-schematic-cpu}.
    \item {\bf Data bus}: This group of signals is used to send the data that the CPU wants to write
    to memory or to an I/O device, or to receive the data that the CPU wants to read from memory or
    an I/O device. We can see the data bus is bidirectional, as the CPU can read or write data from
    the data bus. The data bus is shown as a bidirectional bus in the bottom right of Figure
    \ref{fig:artemisa-schematic-cpu}.
    \item {\bf Control bus}: This group of signals is used to indicate the type of operation that
    the CPU wants to perform. Some of them are inputs used by bus devices to transmit information to
    the CPU, while others are outputs used by the CPU to control the bus devices. The control bus is
    shown in the left side of Figure \ref{fig:artemisa-schematic-cpu}.
\end{itemize}

We will have a closer look at each of these groups in the following sections. But before that, it is
necessary to understand the interface of the Z80 CPU.

\subsection{CPU Interface}

The Z80 CPU has a set of pins that are used to connect it to the rest of the system. These pins are
used to provide power to the CPU and to communicate to other components of the system.

The Z80 CPU has 40 pins. Two of these pins, \sym{VCC} and \sym{GND}, are used to provide power to
the CPU. The rest, used to communicate to the system bus, are described as follows:

\begin{itemize}
    \item \sym{CLK}: This pin is used to provide the clock signal to the CPU. The Z80 CPU requires a
    clock signal to operate. The clock signal is used to synchronize the internal operations of the
    CPU. In the MSX System, the Z80 CPU requires a clock signal with a frequency of 3.58 MHz.    
    \item \sym{RESET}: This pin is used to reset the CPU. When the \sym{RESET} pin is low, the CPU
    will reset and start executing instructions from the address 0x0000. 
    \item \sym{A0} to \sym{A15}: These pins are the Address Bus. They are used to send the address
    of the memory location or the port of the I/O device that the CPU wants to read or write. The
    Z80 CPU has a 16-bit address bus, which means that it can address up to $2^{16}$ memory
    locations, or 64KB of memory. 
    \item \sym{D0} to \sym{D7}: These pins are the Data Bus. They are used to send the data that the
    CPU wants to write to memory or to an I/O device, or to receive the data that the CPU wants to
    read from memory or an I/O device. The Z80 CPU has an 8-bit data bus, which means that it can
    read or write 8 bits of data at a time.
    \item \sym{MREQ} and \sym{IORQ}: These pins are used to indicate if the CPU wants to access the
    memory or an I/O device. When the \sym{MREQ} pin is low, the CPU is accessing the memory. When
    the \sym{IORQ} pin is low, the CPU is accessing an I/O device.
    \item \sym{RD} and \sym{WR}: These pins are used to indicate if the CPU wants to read or write
    data when \sym{MREQ} and \sym{IORQ} are active low. When the \sym{RD} pin is low, the CPU is
    reading data from memory or an I/O device. When the \sym{WR} pin is low, the CPU is writing data
    to memory or an I/O device.
    \item \sym{M1} and \sym{RFSH}: These pins are used to indicate if the CPU is fetching an
    instruction or \toref{refreshing the memory}. When the \sym{M1} pin is low, the CPU is fetching
    an instruction. When the \sym{RFSH} pin is low, the CPU is refreshing the memory.
    \item \sym{WAIT}: This pin is used to indicate if the CPU have to wait for a slow device to
    finish its operation. When the \sym{WAIT} pin is low, the CPU will wait an additional clock
    cycle before continuing the read or write operation. 
    \item \sym{INT} and \sym{NMI}: These pins are used to indicate if the CPU is receiving an
    interrupt from a peripheral. When the \sym{INT} pin is low, the CPU is receiving a maskable
    interrupt. When the \sym{NMI} pin is low, the CPU is receiving a non-maskable interrupt. The
    difference between maskable and non-maskable interrupts is that the CPU can ignore the former
    but not the latter using the \sym{DI} instruction. In MSX System, the usage of non-maskable
    interrupts is not considered, so the \sym{NMI} pin is wired to high to be always inactive.
    \item \sym{HALT}: This pin is used to indicate if the CPU is halted. When the \sym{HALT} pin is
    low, the CPU is halted as result of executing the \sym{HALT} instruction. In this state, the CPU
    stops executing instructions until an interrupt is received. This output signal may be used by
    computer designers to react to the CPU being halted. For example, triggering a reset signal to
    restart the CPU. This pin is not used in the MSX system, so it is unconnected.
    \item \sym{BUSREQ} and \sym{BUSACK}: These pins are used to request the control of the system
    bus by some device other than the CPU. When the \sym{BUSREQ} pin is low, the device is
    requesting the control of the system bus to the CPU. Once the CPU has finished its current
    machine cycle, it will release the control of the system bus by setting the \sym{BUSACK} pin
    low. This pair of pins is not used in the MSX system, so \sym{BUSREQ} is wired to high to be
    always inactive and \sym{BUSACK} is unconnected.
\end{itemize}

\begin{theory}{Memory refreshing}
    The Z80 CPU has a built-in memory refresh circuit that is used to refresh the dynamic memory
    used in the system. 
    
    Dynamic memory (DRAM) is a type of memory that is built using one tiny capacitor to store a bit
    of data in each memory cell. The electric charge of the capacitor determines the value of the
    bit stored in the memory. As each bit cell is made by just a capacitor, they require much less
    space than static memory, which uses many transistors to implement the flip-flop to store a
    single bit. As result, dynamic memory is cheaper and can store more data in the same space than
    static memory.
    
    However, the dynamic memory has a drawback: the charge of the capacitors will leak over time, so
    the data stored in the memory will be lost in a matter of seconds. To prevent this, the data
    stored in the memory must be refreshed periodically by reading the data stored in the memory,
    what makes the internal circuitry of the memory chip to recharge the capacitor if needed. This
    operation is called memory refreshing.
    
    The Z80 CPU has a built-in mechanism to refresh the memory. After each instruction opcode fetch,
    the CPU generates a refresh cycle while it decodes the opcode. During the refresh cycle, the CPU
    activates \sym{MEMRQ}, \sym{RFSH} and \sym{RD} signals and puts the address to be refreshed in
    the address bus. Any device using DRAM attached to the system bus may use this cycle to refresh
    its bit cells. In order to ensure the whole 64KB memory space is refreshed, the Z80 CPU has a
    refresh address counter that is incremented after every refresh cycle. Hence, the whole 64KB
    memory space is refreshed several times per second as program instructions are executed.

    Artemisa uses static memory, so the memory refreshing is not needed for its own RAM. However,
    the refresh signals are connected to the system bus to allow any external device connected to
    the cartridge slots to use the refresh cycle to refresh its memory. Also, this signal is needed
    to distinguish a regular memory access from a refresh cycle, as we will see in the next chapter.
\end{theory}

\subsection{Interface to the address bus}

The 16 address lines of the Z80 CPU are not directly connected to the system bus. Instead, they are
connected to a set of buffers that are used to drive the address bus as shown in Figure
\ref{fig:artemisa-schematic-cpu-address}. These buffers are used to isolate the CPU from the rest of
the system, and to ensure that the CPU can drive the address bus without being affected by the other
components of the system. The buffers used in the design are \sym{74HCT244} 3-state buffers with
\toref{TTL-compatible inputs}. 

\begin{figure}[htbp]
    \centering
    \includegraphics[width=\linewidth,trim={9cm 14.45cm 2cm 8cm},clip]{figures/artemisa-schematic-cpu}
    \caption{Schematic diagram of CPU interface to the address bus}
    \label{fig:artemisa-schematic-cpu-address}
\end{figure}

\begin{theory}{TTL-compatible inputs}
    TTL-compatible inputs in CMOS device are inputs that are compatible with the TTL logic family
    and its voltage levels.

    As introduced in Chapter 3, Artemisa uses CMOS technology in most of its integrated circuits.
    One of the consequences is that the voltage levels used in the digital circuits are different
    from the voltage levels used in the TTL logic family, as shown in Figure
    \ref{fig:voltage-levels}.

    When voltage levels are mixed, we must ensure that a device of one logic family will produce
    outputs in a voltage range the device of the other logic family connected to that output can
    understand. The following combinations of voltage levels are possible between CMOS and TTL:
    
    \begin{itemize}
    \item A CMOS low output voltage is in the range of 0V to 0.05V. As the TTL logic family expects
    from 0V to 0.8V to represent a logical 0, the voltage levels are directly compatible and a TTL
    device can be connected to the CMOS output without any issue.
    \item A CMOS high output voltage is in the range of 4.95V to 5V. As the TTL logic family expects
    from 2V to 5V to represent a logical 1, the voltage levels are directly compatible and a TTL
    device can be connected to the CMOS output without any issue.
    \item A TTL low output voltage is in the range of 0V to 0.5V. As the CMOS logic family expects
    from 0V to 1.5V to represent a logical 0, the voltage levels are directly compatible and a CMOS
    device can be connected to the TTL output without any issue.
    \item A TTL high output voltage is in the range of 2.7V to 5V. As the CMOS logic family expects
    from 3.5V to 5V to represent a logical 1, there is a compatibility issue. We must be prepared
    for a TTL device that produces a high signal with a voltage between 2.7V and 3.5V, which will
    not be recognized as a logical 1 by the CMOS devices connected to it.
    \end{itemize}

    To solve this compatibility issue, the CMOS device must have TTL-compatible inputs if outputs
    from a TTL device are connected to it. The Z80 CPU used in Artemisa is a Z84C00, which uses CMOS
    voltage levels in its outputs. However, if you want to use a spare unit of an NMOS Z8400 CPU,
    you will find it uses TTL output levels instead. Thanks to the TTL-compatible inputs of the
    \sym{74HCT244} buffers, the NMOS version of the Z80 CPU can be connected to the system bus
    without any issue.
\end{theory}

\begin{figure}[htbp]
    \centering
    \includegraphics[width=0.5\linewidth,trim={7cm 1cm 7cm 1cm},clip]{figures/voltage-levels}
    \caption{Voltage level comparison between CMOS and TTL logic families}
    \label{fig:voltage-levels}
\end{figure}

This IC has two groups of 4 buffers each, so two ICs are used to drive the 16 address lines of the
Z80 CPU. Each group of buffers has an enable input signal that is used to enable or disable the
buffers. In the MSX System, the CPU has a full control of the system bus for data transfers, and it
is always driving the address bus. Hence, the enable input signals of the buffers are wired to
\sym{GND} be active low.

\subsection{Interface to the data bus}

As with the address bus, the 8 data lines of the Z80 CPU are not directly connected to the system
bus. Instead, they are connected to a set of buffers that are used to drive the data bus as shown in
Figure \ref{fig:artemisa-schematic-cpu-data}. However, in this case, the buffers are bidirectional,
as the CPU can read or write data from the data bus. 

\begin{figure}[htbp]
    \centering
    \includegraphics[width=\linewidth,trim={9cm 10.5cm 2cm 14.7cm},clip]{figures/artemisa-schematic-cpu}
    \caption{Schematic diagram of CPU interface to the data bus}
    \label{fig:artemisa-schematic-cpu-data}
\end{figure}

The buffers used in the design are \sym{74HCT245} 3-state buffers with TTL-compatible inputs. The \sym{DATADIR} signal coming from the decoding logic is used to control the direction of the buffers. 

\begin{itemize}
\item When the \sym{DATADIR} signal is high, the data is transmitted from A inputs to B inputs. This
is the write mode for the CPU. 
\item When the \sym{DATADIR} signal is low, the data is transmitted from B inputs to A inputs. This
is the read mode for the CPU.
\end{itemize}

As with the address bus, the enable input signals of the buffers are wired to \sym{GND} to be always
active.

\subsection{Interface to the control bus}

The control bus is more heterogeneous than the address and data buses due to the mix of input and
output signals. As in the previous cases, the outputs from the CPU are not directly connected to the
system bus and \sym{74HCT244} buffers are used again to drive the control bus as shown in Figure
\ref{fig:artemisa-schematic-cpu-control}.

\begin{figure}[htbp]
    \centering
    \includegraphics[width=\linewidth,trim={1.2cm 10.5cm 9.5cm 10cm},clip]{figures/artemisa-schematic-cpu}
    \caption{Schematic diagram of CPU interface to the control bus}
    \label{fig:artemisa-schematic-cpu-control}
\end{figure}

However, the inputs signal coming from the system bus into the CPU are not buffered. The reason for
it is that those signals come from other circuits of the same board that are designed to drive the
CPU inputs directly. Also, thse inputs come from CMOS devices so voltage level compatibility issues
do not arise.

Among the input signals of the CPU, the case of the \sym{WAIT} signal is special. This signal is
used to indicate if the CPU have to wait additional clock cycles for a slow device to finish a read
or write operation operation. It comes from a non trivial logic that we will see in the next
section.

\subsection{Generation of wait cycles}

The Z80 CPU has a different way to read from memory when the data to be fetch is an instruction
opcode than a regular data read.

When the CPU is fetching an instruction, it will read the opcode expecting the memory to respond
with the opcode data at the end of the second clock cycle after the memory request signal is
activated. This is called a M1 cycle. In constrast, in a regular read operation the CPU expects the
data to be available at the end of the third clock cycle after the memory request signal is
activated.

At the time when the MSX system was designed, the time constraints of a M1 cycle could not be met by
afordable memory devices available in the market. The solution was to use a wait state generator to
insert an additional clock cycle in the M1 machine cycles. This generator ensures the \sym{WAIT}
signal is active during the second clock cycle of a M1 operation, and deactivated after that.

This wait state generator is implemented using a \sym{74HCT74} with two D-type flip-flops chained as
shown in Figure \ref{fig:artemisa-schematic-cpu-control}. In Figure
\ref{fig:wavedrom-wait-generator}, a wave diagram of the signals involved in the wait state
generator is also shown.

\begin{figure}[htbp]
    \centering
    \includegraphics[width=\linewidth,trim={0cm 25cm 0.5cm 0cm},clip]{figures/wavedrom-wait-generator}
    \caption{Wave diagram of the wait state generator}
    \label{fig:wavedrom-wait-generator}
\end{figure}

The first flip-flop stores the state of the \sym{/M1} signal on the rising edge of the clock signal.
As \sym{/M1} is active-low during the first cycle of a M1 operation (point {\it a} in Figure
\ref{fig:wavedrom-wait-generator}), the output of that flip-flop will be low during the second clock
cycle, activating the \sym{/WAIT} signal of the CPU (point {\it b} in Figure
\ref{fig:wavedrom-wait-generator}) and causing it to wait an additional clock cycle for the memory
to respond. 

On the third clock cycle, the low output coming from the first flip-flop (point {\it c} in Figure
\ref{fig:wavedrom-wait-generator}) will be loaded into the second flip-flop on the rising edge of
the clock (point {\it d} in Figure \ref{fig:wavedrom-wait-generator}). As the output of that
flip-flop is connected to the \sym{/SD} input of the first flip-flop, the third clock cycle will
cause the first flip-flop to be set again to high no matter its input (point {\it e} in Figure
\ref{fig:wavedrom-wait-generator}). Deactivating the \sym{/WAIT} signal of the CPU. 

At the end of this third cycle, the CPU will sample the \sym{/WAIT} signal to find it deactivated.
Then, the CPU reads the opcode from the data bus, deactivates \sym{/M1} and the rest of signals of
the M1 cycle. The CPU will not receive any other \sym{/WAIT} pulse until the next M1 cycle.

Thanks to this small circuit, the computer enforces the generation of its own wait states to ensure
the memory can respond in time to the shorter M1 cycles of the Z80 CPU. However, what if the CPU is
accessing any device that cannot respond in just three clock cycles? 

To cover this case, the system bus carries a \sym{/EXTWAIT} signal from the peripherals that can be
activated by any device that needs more time to respond to a read or write operation. If you look at
Figure \ref{fig:artemisa-schematic-cpu-control}, you will see that the \sym{/EXTWAIT} signal is
connected to the \sym{/RD} input of the first flip-flop of the wait state generator and the
\sym{/SD} input of the second flip-flop. If \sym{/EXTWAIT} is activate-low, the first flip-flop will
be set to low, activating the \sym{/WAIT} input of the CPU, while the second flip-flop will be set
to high, ensuring the first flip-flop is not set to high by it. This way, the \sym{/WAIT} signal is
active during the period the \sym{/EXTWAIT} signal is active. How the \sym{/EXTWAIT} signal is
generated will be explained in Chapter 10.

\section{Circuit assembly}

\subsection{CPU socket}

The first thing to do is to assemble the CPU socket to the motherboard. The CPU socket is a 40-pin
DIP socket that is used to connect the Z80 CPU to the motherboard. The socket is soldered to the top
side of the motherboard, as shown in Figure \ref{fig:mount-cpu-01}.

\begin{figure}[htbp]
    \centering
    \includegraphics[width=0.8\linewidth]{figures/mount-cpu-01}
    \caption{Placement of CPU socket {\tt U6}}
    \label{fig:mount-cpu-01}
\end{figure}

\begin{warning}[Do not connect the CPU to the socket yet!]
    At this point, the input connections to the CPU are still disconnected. The CMOS devices are
    very sensible to static electricity, so it is recommended to handle them with care. To avoid
    damaging the CPU, do not connect it to the socket until the connections are made.
\end{warning}

Once the CPU socket is soldered to the motherboard, the next step is to solder the decoupling
capacitor \sym{C3} of 100nF. You will find its placement on the right side of the CPU socket, as
shown in Figure \ref{fig:mount-cpu-02}.

\begin{figure}[htbp]
    \centering
    \includegraphics[width=0.8\linewidth]{figures/mount-cpu-02}
    \caption{Placement of CPU decoupling capacitor {\tt C3}}
    \label{fig:mount-cpu-02}
\end{figure}

\subsection{Address bus buffers}

Solder the two \sym{74HCT244} ICs \sym{U3} and \sym{U5} to the motherboard. The ICs are placed on
the left side of the CPU as shown in Figure \ref{fig:mount-cpu-03}.

\begin{figure}[htbp]
    \centering
    \includegraphics[width=0.8\linewidth]{figures/mount-cpu-03}
    \caption{Placement of address bus buffers \sym{U3} and \sym{U5}}
    \label{fig:mount-cpu-03}
\end{figure}

Once the ICs are soldered to the motherboard, the next step is to connect their decoupling
capacitors \sym{C15} and \sym{C16} of 100nF. You will find their placement close to the power pin of
the ICs, as shown in Figure \ref{fig:mount-cpu-04}.

\begin{figure}[htbp]
    \centering
    \includegraphics[width=0.8\linewidth]{figures/mount-cpu-04}
    \caption{Placement of buffer decoupling capacitors \sym{C15} and \sym{C16}}
    \label{fig:mount-cpu-04}
\end{figure}

\subsection{Data bus buffers}

Solder the two \sym{74HCT245} ICs \sym{U8} to the motherboard. The IC is placed on the right side of
the CPU as shown in Figure \ref{fig:mount-cpu-05}.

\begin{figure}[htbp]
    \centering
    \includegraphics[width=0.8\linewidth]{figures/mount-cpu-05}
    \caption{Placement of data bus buffers \sym{U8}}
    \label{fig:mount-cpu-05}
\end{figure}

After soldering \sym{U8}, it is time for its decoupling capacitor \sym{C4} of 100nF. It is placed close to the power pin of the IC, as shown in Figure \ref{fig:mount-cpu-06}.

\begin{figure}[htbp]
    \centering
    \includegraphics[width=0.8\linewidth]{figures/mount-cpu-06}
    \caption{Placement of buffer decoupling capacitor \sym{C4}}
    \label{fig:mount-cpu-06}
\end{figure}

\subsection{Control bus buffers}

Solder the \sym{74HCT244} IC \sym{U4} to the motherboard. The IC is placed on the right side of the
CPU, just below the data bus buffers, as shown in Figure \ref{fig:mount-cpu-07}.

\begin{figure}[htbp]
    \centering
    \includegraphics[width=0.8\linewidth]{figures/mount-cpu-07}
    \caption{Placement of control bus buffer \sym{U4}}
    \label{fig:mount-cpu-07}
\end{figure}

After soldering \sym{U4}, it is time for its decoupling capacitor \sym{C17} of 100nF. It is placed close to the power pin of the IC, as shown in Figure \ref{fig:mount-cpu-08}.

\begin{figure}[htbp]
    \centering
    \includegraphics[width=0.8\linewidth]{figures/mount-cpu-08}
    \caption{Placement of buffer decoupling capacitor \sym{C17}}
    \label{fig:mount-cpu-08}
\end{figure}

\subsection{Wait state generator}

Solder the \sym{74HCT74} IC \sym{U7} to the motherboard. The IC is placed on the right side of the
control bus buffers just soldered, as shown in Figure \ref{fig:mount-cpu-09}.

\begin{figure}[htbp]
    \centering
    \includegraphics[width=0.8\linewidth]{figures/mount-cpu-09}
    \caption{Placement of wait state generator \sym{U7}}
    \label{fig:mount-cpu-09}
\end{figure}

After soldering \sym{U7}, it is time for its decoupling capacitor \sym{C19} of 100nF. It is placed close to the power pin of the IC, as shown in Figure \ref{fig:mount-cpu-10}.

\begin{figure}[htbp]
    \centering
    \includegraphics[width=0.8\linewidth]{figures/mount-cpu-10}
    \caption{Placement of wait state generator decoupling capacitor \sym{C19}}
    \label{fig:mount-cpu-10}
\end{figure}

\subsection{Testing the circuit}

In this particular case, there is not much to be tested. Even if we connect the CPU to the socket,
the rest of the devices at the other end of the buses are not connected yet. If you'd power on the
motherboard at this point, the CPU would not be able to communicate with anything. What we actually
have is a working CPU that is alone in the world.

\section{Summary}

In this chapter, we have seen how the CPU is connected to the rest of the system through the system
bus. We have seen how the CPU is connected to the address, data, and control buses, and how the CPU
communicates with the rest of the system through these buses. We have also seen how the CPU is
connected to the system bus through a set of buffers that are used to isolate the CPU from the rest
of the system, and to ensure that the CPU can drive the system bus without being affected by the
other components of the system. Finally, we have assembled the CPU and the buffers to the
motherboard, including the flip-flops used to generate wait states.

In the next chapter, we will see how the output signals from the CPU can be decoded to control the
built-in devices attached to the system bus, such as the memory, the PPI, the PSG and the video
card.